\input{preamble}
\renewcommand{\draftstatus}{\begin{center}\small
  \textbf{Research draft revised on 11 November 2026.}\\
  Completed campaign: 600 scenarios, 16,800 primary trials and 1,680 extra repetitions.
  Descriptive comparisons are supported by the accepted final records.
  Causal component attribution, matched EBM ablation and generalization remain untested.
\end{center}}
\date{11 November 2026}

\title{Which Operational Service Do Energy-Optimization Methods Deliver?\\
Available Time, Economic Quality, Unserved Energy and Generalization}
\begin{document}
\maketitle
\draftstatus
\begin{abstract}
The value of an energy-optimization method depends on the operational service
requested, the scenario family and the available time. We compare direct
HiGHS, rolling-horizon construction with repair (RH), one-candidate JEPA and
a sixteen-candidate energy-based pipeline on 600 scenarios from six families,
with seven independent execution budgets. All 16,800 primary trials and 1,680
additional repetitions are complete. At 120 seconds, HiGHS returns a valid
output for all 600 scenarios and reaches five-percent reference-relative
quality on 505, whereas RH preserves service without unserved energy on 577
scenarios versus 482 for HiGHS. Requiring both five-percent quality and no
unserved energy yields 475 successes for HiGHS, 248 for the pipeline, 222 for
JEPA and 76 for RH. These aggregate rankings conceal targeted compromises.
At ten seconds on VRE-with-batteries cases, the four-method coverage ceiling
for the joint requirement is 59/100 versus 46/100 for HiGHS, including six
cases served only by JEPA. At longer allocations, JEPA can reduce completion
time on selected acceptable cases; the pipeline improves some JEPA quality
counts but rarely supplies the fastest acceptable output. The comparison
includes online embeddings and full process exit, excludes static preload,
and retains failed trials in coverage. Descriptive complementarity is
established; pre-solve selection, independent transfer and an incremental
benefit of energy-guided sampling remain to be tested.

\end{abstract}

\section{Introduction}
An operator may need a dispatch within thirty seconds without unserved energy,
while another decision may allow more time in exchange for a tighter economic
target. These requests need not favor the same optimization method. A
mathematically valid dispatch can satisfy the model while shedding load:
economic penalties make this action costly, but do not prohibit it. We
therefore distinguish physical validity, service continuity and economic
quality, and ask what a method delivers before the decision deadline.

The relevant time is complete execution, including scenario processing,
candidate generation, repair, validation, output writing and process closure.
An incumbent inside a still-running process is not an available operational
response. Conversely, an output close to a common reference may be useful
without constituting a proof of optimality. The operational question is thus
not simply which algorithm has the lowest average objective, but which
combination of time, cost and service it can deliver for a stated need.

Learning to guide mixed-integer optimization has several established forms.
Graph-based branching learns variable-selection decisions within an exact
search procedure \citep{gasse2019}. Neural Diving generates partial assignments
that a solver completes \citep{nair2021}; learned large-neighborhood search
chooses neighborhoods for subsequent optimization \citep{sonnerat2022}.
Our system generates zonal strategy tensors and reconstructs a dispatch
through a shared repair cascade. The paper does not claim the general
invention of learned optimization. It characterizes the usefulness and limits
of these particular constructions alongside direct optimization and RH.

Three questions structure the study. First, how do the available trade-offs
between execution time, economic cost and unserved energy change with scenario
family and allocated budget? Second, where do learned proposals provide a
useful compromise, and which apparent gains can instead be explained by
candidate count or solver repair? Third, would any observed usefulness survive
structural or temporal changes? The completed campaign answers the first
question descriptively and motivates the others; it does not answer the
unexecuted ablation and generalization questions.

\section{Evidence scope and common protocol}
The final campaign contains \textbf{600 scenarios, 100 per family}, all four
methods and all seven budgets: 16,800 primary executions. Two extra
repetitions on 30 preselected scenarios add 1,680 executions; all 18,480
expected records are present. The final acceptance report records 600 valid
references and no invalid or late output incorrectly counted as successful.
Timeouts, absent incumbents and a late exit remain observed failures.

The present tables use primary repetitions, with repeated observations
reported separately. The six families are three criticality strata and three
specialized generator profiles, not six ordered severity levels. References
are certified optimal on 383/600 scenarios; the remainder are validated
incumbents. The final descriptive analysis was computed before this writing
revision and is archived with source hashes in
\path{notes/final_results_2026-11-11.json}. The earlier 126-case narrative is
historical context only. Service profiles were motivated by inspected results,
so the completed cohort is an exploratory development resource for routing,
not an untouched test of the proposed selection rules.

The common physical protocol below retains the historical five-percent
economic endpoint. We additionally assess unserved energy explicitly.

\input{common_protocol}

\section{Operational requirements and reported outcomes}
Let $T_{smb}$ be complete exit time and let $C_{smb}$ be the validated model
objective. This cost includes the unserved-energy penalty; it is not a
fuel-only or service-independent operating cost. For a returned valid dispatch,
define delivered unserved energy in MWh as
\begin{equation}
 E_{smb}=\sum_{z,t}P^{\rm uns}_{smbzt}\Delta t .
 \label{eq:ens}
\end{equation}
The objective and $E_{smb}$ must be reported separately: an economically
acceptable solution can still violate an operator's service requirement.

Write $Q_{smb}(\tau)$ for the cost condition in Eq.~\eqref{eq:economic},
including its absolute-tolerance convention near a zero reference. With
$\tau=\infty$, the cost constraint is disabled, so that $Q_{smb}(\infty)=1$
does not require a reference. For a maximum acceptable unserved energy
$\bar E$, the joint requirement outcome is
\begin{equation}
 Y_{smb}(\tau,\bar E)
 =F_{smb}\,Q_{smb}(\tau)\,\ind\{E_{smb}\leq\bar E\}.
 \label{eq:service}
\end{equation}
An observed execution failure has outcome zero without assigning a fictitious
cost or unserved-energy value. Missing trials or missing information needed
to assess an otherwise returned output remain unknown. On complete observed records, setting
$\bar E=\infty$ recovers the cost-only endpoint; continuity-sensitive requests
instead impose an explicit energy cap. An inactive condition does not require
its underlying measurement.

A zero-unserved-energy profile uses $\bar E=\varepsilon_E=10^{-3}$ MWh,
matching the accepted campaign analysis. This aggregated energy tolerance is
distinct from the power-balance tolerance. Sensitivity to less restrictive
caps is reported explicitly; changing a cap changes the service requested.

A requirement is described by a deadline, a cost tolerance, an unserved-energy
cap and the operator's priority. For example, \emph{minimize cost within thirty
seconds while preserving service} has no mandatory relative-cost threshold.
In contrast, \emph{reach one-percent economic quality as quickly as possible}
places a hard condition on quality and optimizes completion time; service
continuity is required only if an energy cap is also stated. A tolerance to
the risk of failing these requirements is an additional input to the companion
router. These profiles are motivated by already inspected results and are
exploratory on this completed cohort; confirmatory profiles must be frozen before an
independent evaluation.

Coverage counts all observed trials, including failures. Time, cost and energy
summaries on returned solutions additionally state their support. A
retrospective comparison can identify non-dominated outputs at a fixed budget,
but it cannot select them before solving. Nor can a run allocated 120 seconds
that finishes quickly be substituted for an independent thirty-second run:
allocation changes the execution and repair allowances. The existing grid is
not an anytime trajectory or a time-to-target curve.

The amended joint-analysis plan focuses on five distinct deadlines:
$5,10,30,60,120$ seconds. At each deadline, service profiles, coverage and
paired cost/time/ENS contrasts are evaluated separately. Historical one- and
two-second records remain archived; they are not pooled into this amended
service analysis. This scope change does not revise the historical results
reported below and remains exploratory on the inspected development cohort.

\section{Methods and reusable learned representations}
\subsection{Direct solver and rolling-horizon baseline}
The direct HiGHS method constructs the corrected MILP from the scenario and
uses sixteen CPU threads. Its internal time allowance is bounded by the
remaining execution budget after the unchanged finalization-reserve policy.
The reserve is part of the frozen protocol; it is not retuned on evaluation
outcomes. Model construction and extraction remain in the measured interval.

The rolling-horizon method constructs a discrete proposal from scenario physics
and submits it to the same HiGHS repair cascade as the learned methods. The
label ``rolling horizon'' therefore describes its proposal mechanism, not a
guarantee that the final dispatch is heuristic-only or avoids mixed-integer
optimization. All four methods use the same physical acceptance criterion.

\subsection{Scenario graph, autoencoder and latent predictor}
Following the existing methodological manuscript, the scenario is represented
as a heterogeneous temporal graph connecting system, region, zone and equipment
levels. A hierarchical temporal encoder produces scenario embeddings online.
The frozen feature contract and training normalization accompany the encoder;
evaluation scenarios and their embeddings are not preloaded.

Each reference strategy is represented by
\begin{equation}
 u\in\{0,1\}^{Z\times T\times8},\quad
 u_{zt}=[b^{\rm ch},b^{\rm dis},p^{\rm ch},p^{\rm dis},d^{\rm DR},v^{\rm th},u^{\rm th},i^{\rm imp}]_{zt}.
\end{equation}
The storage channels encode positive charging/discharging activity according
to the label extractor, and permit an idle state. The import-mode label is a
system-level binary broadcast over zones. The common repair interface consumes
the first seven channels and retains import/export as a repair-model decision.
Consequently, eight output channels do not imply eight independent operational
degrees of freedom at every location and time.

The autoencoder compresses the strategy into a fixed set of latent tokens and
decodes per-zone, per-time logits. A scenario-to-latent predictor supplies an
anchor without observing the reference strategy. The current AE checkpoint is
configured for at most 300 zones and 24 time steps, with 32 tokens of dimension
192. Its decoder contains learned zone and time index embeddings. The training
support must be audited empirically; architectural capacity is not evidence
that every compatible size was represented in training.
The encoder reads zone trajectories with a BiGRU over $T$, mean-pools in time,
and lets learned query tokens cross-attend over the masked zones. The decoder
builds learned zone--time queries that cross-attend over the strategy tokens,
with a projection of the scenario embedding as a local conditioning signal.
This methodological description is adapted from the existing manuscript;
its previous numerical conclusions are not reused here.

\subsection{Energy sampling and shared repair}
The JEPA method decodes one anchor. The pipeline generates sixteen latent
candidates by seeded Langevin updates using the learned energy, decodes them,
and attempts repair. The accepted result is the least-cost independently valid
candidate completed within the allowed computation interval. Candidates still
running cannot replace that result during finalization. Every child process
must nevertheless terminate before the global deadline.

The first repair stage fixes supplied targets, and later stages may release
selected decisions under deviation and flip penalties. The amended runtime
restores the original integer domains and structural fixes; this cascade may
solve MILPs. Calling it a pure LP reconstruction would misstate the current
algorithm. One-candidate methods allocate sixteen solver threads to their
repair; sixteen-candidate methods allocate one thread to each candidate.

The new service profiles assess the dispatch returned by each unchanged
method. In particular, the pipeline still selects by cost internally; we do
not retrospectively replace its output with a different candidate chosen
using the new energy cap. Such a change would define a new method variant.

\subsection{A planned matched test of the energy contribution}
This control remains a prepared, unexecuted comparison in this revision.
JEPA with $K=1$ versus the pipeline with $K=16$ conflates candidate count and
energy guidance. The isolated \texttt{perturb16} control loads the same models,
uses the same candidate seeds, initialization, noise, drift, clipping, iteration
count and repair allocation, but sets the energy-gradient term to zero. Both
methods receive the same wall-clock deadline and hardware envelope. Their
realized computation times need not be identical because gradient evaluation
has a cost. A difference can therefore be interpreted as the operational value
of the energy term within this specified sampler and resource budget, rather
than as a universal statement about EBMs.

\section{Observed operational trade-offs on the complete cohort}
\subsection{Changing the service changes the ranking}
Table~\ref{tab:final_aggregate} counts every primary scenario, including failed
executions. At 120 seconds HiGHS leads economic coverage whereas RH leads
coverage without ENS. JEPA returns more valid outputs than the pipeline,
but the pipeline satisfies the tighter joint requirement on more scenarios.
These statements concern this implementation and resource envelope.

\begin{table}[htbp]\centering\small
\caption{Completed primary campaign at 120 seconds. Counts out of 600.
Economic quality means $g\leq0.05$; all service and quality counts require
a valid output and full exit within the deadline.}
\label{tab:final_aggregate}
\begin{tabular}{lrrrr}\toprule
Criterion & HiGHS & Pipeline & JEPA & RH\\\midrule
Valid timely output & 600 & 311 & 467 & 577\\
Without ENS & 482 & 265 & 253 & 577\\
Five-percent economic quality & 505 & 269 & 259 & 76\\
Quality and no ENS & 475 & 248 & 222 & 76\\\bottomrule
\end{tabular}
\end{table}

\asset{figures/final_service_coverage.png}
{Final coverage at 30 and 120 seconds, with 100 scenarios in every family.
Columns distinguish validity, continuity and their joint economic requirement.
Budgets correspond to independent runs.}{fig:final_service}

\subsection{Family-specific operational uses}
\paragraph{Low criticality.}
At five seconds HiGHS already returns 100/100 outputs within five percent of
the reference, 99 without ENS, with a median exit time of about 2.1 seconds.
The learned methods can solve this family, but a strong direct baseline leaves
little room for useful acceleration. The single strict-ENS exception is
examined in the threshold sensitivity below.

\paragraph{Medium criticality.}
At ten seconds RH returns 100/100 outputs without ENS versus 25/100 for HiGHS,
while RH's median reference gap is approximately 32.5\% among its successful
outputs. At sixty seconds HiGHS reaches the joint five-percent/no-ENS target
on all 100 cases. At 120 seconds JEPA offers 25 cases combining a faster
exit than HiGHS, five-percent quality and no ENS. The pipeline meets that
joint target on 51 cases versus 36 for JEPA, but returns fewer valid outputs
(79 versus 92) and incurs substantial execution overhead.

\paragraph{High criticality.}
At thirty seconds RH preserves service on 100/100 cases versus 8/100 for HiGHS.
At 120 seconds HiGHS reaches five-percent economic quality on 77 cases and
the joint no-ENS requirement on 61. RH remains fast, with a successful-output
median exit of approximately 11 seconds, but its median reference gap is
about 63.8\%. The joint target is reached by only two pipeline and one JEPA
outputs. The dominant trade-off here is between rapid RH continuity and
stronger HiGHS economic quality at a longer allocation.

\paragraph{VRE with batteries.}
At ten seconds the joint target is reached on 46 cases by HiGHS, 38 by JEPA,
26 by RH and none by the pipeline. Their union covers 59 cases; six are
covered only by JEPA. This is observed complementarity under an identical
deadline and requirement. At 120 seconds HiGHS covers all 100 cases under
the joint target, compared with 80 for the pipeline and 78 for JEPA.
JEPA nevertheless returns a faster, joint-acceptable output than HiGHS on
43 cases. Larger nominal allocations do not establish the same outcome
under a shorter deadline.

\paragraph{Network.}
At thirty seconds RH returns 98 outputs without ENS versus 40 for HiGHS.
At 120 seconds HiGHS satisfies the joint economic/service target on 95 cases.
JEPA supplies 96 valid outputs but only 18 satisfy that joint target;
15 of these are faster than HiGHS. Validity alone would therefore substantially
overstate JEPA's coverage for an operator who also requires continuity.

\paragraph{Mathematician.}
At thirty seconds RH returns 96 outputs without ENS versus one for HiGHS.
Among 95 cases where both return valid outputs, RH has lower cost on 78;
cost differences greater than one euro define this comparison. At 120 seconds
HiGHS improves economic quality but returns only 26 outputs without ENS,
20 meeting the joint five-percent target. RH preserves its 96 no-ENS outputs,
with a median reference gap near 99\%. Neither learned method meets the
joint target in this family at the tested allocations.

\subsection{What the learned methods add}
Table~\ref{tab:jepa_final} reports paired speed opportunities at the same
120-second allocation. Its joint column requires JEPA to satisfy quality and
continuity as well as finish earlier; it is not a statement that all methods
meet the joint requirement on the paired support. Differences of at least
one second and ten percent are shown as a descriptive practical threshold,
not a statistical significance test.
\begin{table}[htbp]\centering\small
\caption{JEPA versus HiGHS at 120 seconds; 100 evaluated scenarios per family.}
\label{tab:jepa_final}
\begin{tabular}{lrrrr}\toprule
Family & Valid pairs & Faster JEPA & Faster + joint & Practical gain + joint\\\midrule
Medium & 92 & 57 & 25 & 18\\
VRE/batteries & 81 & 44 & 43 & 27\\
Network & 96 & 82 & 15 & 14\\\bottomrule
\end{tabular}
\end{table}

The pipeline improves one-percent economic coverage relative to JEPA on
medium-criticality cases (42 versus 27 out of 100) and VRE-with-batteries
cases (72 versus 62) at 120 seconds. HiGHS reaches 100 in both families.
Across the 311 pipeline/HiGHS valid pairs, the pipeline is faster on 17;
only seven also meet the joint five-percent/no-ENS target. It is the fastest
joint-acceptable member on three cases in the full retrospective portfolio,
including a case where HiGHS violates the strict ENS cap. These rare cases
are not a broad winning regime, and neither candidate-count nor energy-model
causality follows from the four-method comparison.

\subsection{Repetitions and sensitivity of the service definition}
Three executions are available for each of 30 size-spanning cases, five per
family. At 120 seconds, all three HiGHS executions return valid outputs on
all 30 cases, but joint quality/no-ENS acceptance changes between repetitions
on six. The corresponding joint-acceptance variability counts are three for
the pipeline, one for JEPA and zero for RH. RH's stability accompanies limited
economic coverage; it is not evidence of higher quality. Repetitions are not
additional independent scenarios, and this subset does not establish
deployment risk guarantees.

At 120 seconds a perfect retrospective selector covers 477 scenarios under
the joint five-percent/no-ENS criterion versus 475 for HiGHS. The two additions
are \path{low_criticality_scenarios/scenario_00066} and
\path{persona_network_operator_scenarios/scenario_00039}; HiGHS returns about
0.121 and 0.682 MWh of ENS respectively. With an allowable cap of 1 MWh,
HiGHS and the portfolio both cover 477 cases. The small coverage gain is thus
specific to a strict service request, not robust evidence of a large general
advantage. Energy caps must have an operational interpretation.

\subsection{A bounded opportunity for routing}
For continuity at thirty seconds, the best fixed method covers 575/600 cases,
the best retrospective family-conditioned fixed rule 593, and a perfect
individual choice 597. For cost-only five-percent quality at 120 seconds,
HiGHS and the portfolio both cover 505. At ten seconds, the joint
quality/no-ENS ceiling is 182 versus 168 for the best fixed method.
The opportunity therefore depends on the requested deadline and service.

Among the 475 cases where HiGHS itself meets the joint target at 120 seconds,
the perfect fastest acceptable choice reduces mean exit time from 36.25 to
34.16 seconds, approximately 5.8\%, before any selection overhead. On 68 cases
it saves at least one second and ten percent; their median saving is 5.9
seconds. These are retrospective ceilings on realized outcomes, not measured
router performance, expected deployment gains or an optimality guarantee.
The companion routing paper distinguishes this evidence from all prospective
prediction and online-evaluation claims.

\section{Planned robustness and generalization tests}
The designs below remain prospective. They do not contribute observed results
to the completed 600-scenario map, and no data generation, training or experiment is
launched in this revision. Their purpose is to test whether the relevant
time--cost--service profiles would survive changes in the input distribution.
\subsection{Stage A: frozen-model structural evaluation}
The planned new cases would use independent identifiers and deterministic seeds. Separate
planned development and final cohorts would each contain twenty cases per family. The
generator uses the family parameter spaces, with a documented draw rule that
does not reproduce the historical outcome-independent criticality ranking.
This is a new sampled distribution, not an unqualified IID replication.

The planned controlled suite would use five bases per family, giving thirty cases per
condition, and keeps a distinct development suite. Each base has a 24-hour
control, a doubled-system variant within 300 zones, a line-rewiring variant,
a fleet-capacity variant and a consistent zone-renumbering variant. Renumbering
maps demands, capacities, line endpoints and the external-trade anchor together;
it should preserve the physical optimization problem up to relabeling. It
directly probes sensitivity to learned index embeddings. Capacity changes keep
initial and terminal storage quantities consistent. Doubled cases change both
scale and the linking structure, so their effect is not attributed to size
alone without further controls.

The dimension audit would report actual training shapes, scenario descriptors and
free model integers. ``Extrapolation'' is reserved for dimensions or structures
outside the observed training support; a new draw inside the generator support
is interpolation. Standard solvers already construct a model for new topologies
without training. Learned generalization is useful only if it preserves a
favorable time--cost--service profile against those solvers.

\subsection{Stage B: a distinct variable-horizon model}
The existing 24-step decoder is not applied beyond its capacity. A proposed separate
architecture would replace learned index queries with scenario-conditioned queries
and dynamically generated sinusoidal time positions. The AE, latent predictor
and energy network are initialized and trained anew in that separate planned study on 12-, 24- and 48-hour
cases; 72-hour cases are reserved for extrapolation. The frozen HTE weights
and normalization are reused with an explicitly extended sinusoidal time buffer.
This is partial reuse of a representation, not zero-shot transfer of the old
UC stack.

The prepared corpus design calls for 180 independent bases, each represented at three
training horizons. All horizons of a base share the same training/validation
assignment. Batch size one with gradient accumulation handles variable shapes
without padded observations. Label validity and certification are stored
separately. The prepared AE loss is weighted binary cross-entropy; the predictor
combines latent regression and decoded binary supervision; the energy model
learns a contrastive ranking of reference latents and local noisy proposals.
These choices define an executable new baseline, not a claim to reproduce the
historical training losses exactly.

The horizon controls preserve the first day and extend it using specified daily
scalings of demand and renewable availability, while maintaining hourly
resolution and terminal constraints. Their structured synthetic nature limits
claims about weather or demand forecast shift. Training and final structural
test bases remain disjoint.

\subsection{External solver comparison}
The planned external comparison would evaluate Gurobi directly on 120 existing cases, twenty per family, selected
by a fixed hash before its outcomes are observed. Its repair backend is not
inserted into the four-method campaign. Its prerequisite and license checks
must pass before the extension environment is frozen. Comparisons use common
available cases and the same HiGHS-based reference cost; a negative Gurobi gap
is retained as evidence of reference improvement.

\section{Analyses that remain pending}
The final coverage counts, paired contrasts, repetition summaries and
retrospective service-profile ceilings have been reconciled with the accepted
row-level analysis. They no longer await extraction. What remains is grouped
uncertainty estimation, sensitivity to relative and absolute economic targets,
and a decomposition of ENS and cost by physical scale. These analyses must
retain all observed failures and distinguish conditional successful supports.

\pending{Populate stage-time and repair-change analyses, the matched
\texttt{perturb16} comparison, and structural, variable-horizon and
external-solver evaluations. None is inferred from the completed primary
comparison. No additional experiment is run for this manuscript revision.}
Candidate stages may overlap, so elapsed components cannot be summed without
respecting their execution structure. A valid dispatch after extensive repair
does not by itself demonstrate an accurate learned proposal. Likewise, the
pipeline--JEPA difference cannot isolate the energy-gradient contribution.

\pending{The companion criticality study must test whether pre-solve
components explain incremental outcome variation beyond size and family.
The routing study must establish predictive recovery and timed benefit on
independent data. Family-average contrasts are not fitted routing rules.}

\section{Discussion and limitations}
The completed evidence supports conditional operational roles. RH is useful
when rapid continuity dominates precision, HiGHS when economic accuracy or
already-easy cases dominate, and JEPA on selected time--cost--service
compromises. The pipeline can improve JEPA's economic coverage while remaining
less attractive than direct optimization for most measured requests. No method
is assigned an artificial winning domain.

ENS is penalized inside the cost but remains an independent service outcome.
The economic objective includes export revenues, so net costs can be negative;
relative gaps use the absolute reference denominator. Absolute currency
differences and energy quantities should accompany percentage comparisons.
Near-reference costs are not optimality certificates, and a zero-ENS profile
is not a security, contingency or forecast-robustness guarantee.

The profiles and narrative were developed while the cohort was being explored.
Their final counts describe this cohort, not a preregistered confirmatory test
or calibrated deployment probabilities. Repetitions cover 30 scenarios rather
than the entire cohort. Synthetic families, a capacitated transport network,
aggregate zonal commitment, frozen learned models, one machine and one shared
HiGHS repair backend limit external validity. Low and high families were
selected using criticality and related generator controls; their separation
does not validate the index independently of size or family.

Component effects, useful index reweighting, generalized topologies and a
deployable router require distinct evidence. The companion papers preserve
this distinction between arithmetic index sensitivity, out-of-sample
performance prediction and actual decision benefit.

\section{Conclusion}
Changing the requested service changes the useful method. The final campaign
demonstrates rapid continuity from RH on difficult families, strong economic
coverage from HiGHS, and targeted JEPA gains including unique joint-service
coverage on VRE-with-batteries cases at ten seconds. The pipeline improves
some JEPA quality counts but its execution overhead limits its current role.

Routing potential is request-dependent: coverage headroom is appreciable for
some short-budget requests, small after a family-conditioned baseline for
continuity at thirty seconds, and absent for cost-only five-percent quality
at 120 seconds. Time gains remain possible among already acceptable outputs.
Predicting those gains before solving, at sufficient reliability and after
selection overhead, is the prospective claim of the routing paper.

\section*{Data, provenance and execution status}
\input{notes/final_campaign_provenance}

\bibliographystyle{unsrtnat}\bibliography{paper2/references}
\end{document}
